
\documentclass[aspectratio=169, 10pt]{beamer}

% --- Paquetes Esenciales ---
\usepackage[utf8]{inputenc}
\usepackage[spanish]{babel}
\usepackage{amsmath, amsfonts, amssymb}
\usepackage{booktabs}
\usepackage{graphicx}
\usepackage{tikz}
\usepackage{hyperref}

% --- Configuración de Estilo y Tema Moderno ---
\usetheme{Madrid}
\useinnertheme{rectangles}

% Paleta de Colores Corporativa / Académica
\definecolor{NavyDark}{RGB}{15, 34, 64}       % Azul marino profundo
\definecolor{SlateBlue}{RGB}{41, 74, 110}     % Azul pizarra intermedio
\definecolor{AccentTeal}{RGB}{20, 140, 130}   % Acento verdoso/teal
\definecolor{SoftGray}{RGB}{245, 247, 250}    % Fondo claro
\definecolor{Charcoal}{RGB}{40, 40, 40}       % Texto principal
\definecolor{AlertRed}{RGB}{190, 40, 40}      % Para destacar el error empírico

\setbeamercolor{palette primary}{bg=NavyDark, fg=white}
\setbeamercolor{palette secondary}{bg=SlateBlue, fg=white}
\setbeamercolor{palette tertiary}{bg=NavyDark, fg=white}
\setbeamercolor{structure}{fg=NavyDark}
\setbeamercolor{frametitle}{bg=NavyDark, fg=white}
\setbeamercolor{title}{bg=NavyDark, fg=white}
\setbeamercolor{block title}{bg=SlateBlue, fg=white}
\setbeamercolor{block body}{bg=SoftGray, fg=Charcoal}
\setbeamercolor{block title alerted}{bg=AlertRed, fg=white}
\setbeamercolor{block body alerted}{bg=SoftGray, fg=Charcoal}

\setbeamertemplate{navigation symbols}{}
\setbeamertemplate{footline}[frame number]

% --- Metadatos de la Presentación ---
\title[Clustering Espectral en el S\&P 500]{\textbf{¿Reflejan los Sectores de Wall Street el Riesgo Real?}}
\subtitle{Clustering Espectral y Detección de Comunidades en Redes Complejas}
\author[Arias \& Moreno]{\textbf{Miguel Ignacio Arias Meriño} \and \textbf{Jorge Moreno}}
\institute[UAI]{
	Álgebra Lineal y Optimización para Data Science\\
	\vspace{0.2cm}
	\textit{Red de Correlación Bursátil del S\&P 500 (2014--2024)}
}
\date{Octubre 2026}

\begin{document}
	
	% =============================================================================
	% SLIDE 1: PORTADA
	% =============================================================================
	\begin{frame}
		\titlepage
		\note{
			[0:00 - 0:35] Miguel:
			Buenos días profesores y compañeros. Junto a Jorge Moreno presentaremos
			nuestro trabajo: Clustering espectral y detección de comunidades en redes complejas.
			El riesgo financiero no viaja en línea recta: viaja a través de una red densa y ruidosa.
			Les mostraremos cómo el álgebra lineal nos permite resolver un problema computacionalmente
			intratable y cómo el espectro revela la verdadera estructura del mercado sin supervisión.
		}
	\end{frame}
	
	% =============================================================================
	% SLIDE 2: EL CONFLICTO Y PREGUNTA DE INVESTIGACIÓN
	% =============================================================================
	\begin{frame}{El Enigma: ¿Etiquetas Burocráticas o Riesgo Sistémico?}
		\begin{columns}[c]
			\begin{column}{0.50\textwidth}
				\begin{block}{Taxonomía GICS Tradicional}
					\begin{itemize}
						\item Clasificación estática basada en productos y balances contables.
						\item Cuatro sectores bajo estudio: \textbf{Tecnología, Finanzas, Salud y Energía}.
						\item \textbf{Premisa:} Las empresas de un mismo sector co-varían armónicamente.
					\end{itemize}
				\end{block}
				\vspace{0.2cm}
				\begin{alertblock}{La Falla en Escenarios de Crisis}
					El contagio de liquidez y volatilidad ignora las etiquetas administrativas.
				\end{alertblock}
			\end{column}
			
			\begin{column}{0.48\textwidth}
				\begin{block}{Pregunta Central de Investigación}
					\textit{¿Es capaz el clustering espectral de reconstruir a ciegas las comunidades estructurales de riesgo en el S\&P 500, basándose puramente en fluctuaciones de precios y sin etiquetas previas?}
				\end{block}
				\vspace{0.2cm}
				\textbf{¿Por qué el enfoque espectral?}
				\begin{itemize}
					\item Resuelve la maldición de la dimensionalidad.
					\item Filtra el ruido estocástico del mercado bursátil.
				\end{itemize}
			\end{column}
		\end{columns}
		\note{
			[0:35 - 1:45] Miguel:
			Todo portafolio tradicional se construye bajo la taxonomía GICS. A una empresa se le
			asigna una etiqueta contable: Tecnología, Finanzas, Salud o Energía. Sin embargo, en una crisis,
			los precios no respetan las etiquetas administrativas: responden a dinámicas de covarianza.
			De aquí nace la pregunta: ¿logra el clustering espectral reconstruir estas comunidades a ciegas?
			Es la herramienta idónea porque los autovalores y autovectores actúan como un filtro de ruido topológico.
		}
	\end{frame}
	
	% =============================================================================
	% SLIDE 3: PREPARACIÓN Y TOPOLOGÍA DE LA RED
	% =============================================================================
	\begin{frame}{Construcción y Preparación del Grafo Bursátil}
		\begin{columns}[c]
			\begin{column}{0.52\textwidth}
				\textbf{Pipeline Metodológico:}
				\begin{enumerate}
					\item \textbf{Series Temporales:} 10 años de cotizaciones diarias (2014--2024) vía Yahoo Finance (\(T = 2.515\) ruedas).
					\item \textbf{Universo de Inversión:} 160 activos iniciales; filtrado estricto de NaNs y deslistes \(\to\) \textbf{\(n = 141\) nodos válidos}.
					\item \textbf{Pesos (\(W\)):} Retornos porcentuales diarios y matriz de correlación de Pearson \(\rho_{ij} \in [-1, 1]\).
				\end{enumerate}
				\vspace{0.2cm}
				\textbf{Decisiones Topológicas Clave:}
				\begin{itemize}
					\item \textbf{Sin auto-bucles:} \(\text{diag}(A) = 0\) (evita redundancia estructural).
					\item \textbf{Umbral de similitud:} \(A_{ij} = \rho_{ij}\) si \(\rho_{ij} \ge 0.40\), sino $0$.
				\end{itemize}
			\end{column}
			
			\begin{column}{0.45\textwidth}
				\begin{block}{Diagnóstico Estructural de la Red}
					\begin{itemize}
						\item \textbf{Nodos finales (\(n\)):} 141
						\item \textbf{Aristas ponderadas (\(|E|\)):} 1.984
						\item \textbf{Grado promedio (\(\bar{d}\)):} 28,14 conexiones/nodo
						\item \textbf{Componentes conexas (\(c\)):} 1 (Grafo conexo)
						\item \textbf{Nodos aislados:} 0
					\end{itemize}
				\end{block}
				\vspace{0.1cm}
				\small\textit{Garantiza una matriz de adyacencia no negativa simétrica \(A = A^\top\) libre de componentes triviales.}
			\end{column}
		\end{columns}
		\note{
			[1:45 - 3:15] Miguel:
			Descargamos una década de historia diaria desde Yahoo Finance (2014 a 2024, 2515 ruedas).
			Partimos con 160 activos y descartamos 12 por datos incompletos, fijando 141 nodos limpios.
			Calculamos la correlación de Pearson y tomamos tres decisiones: eliminamos auto-bucles poniendo
			la diagonal a cero; filtramos ruido forzando a cero correlaciones menores a 0.40 y negativas;
			y verificamos que el grafo resultante fuera conexo, con un promedio de 28 aristas por nodo.
		}
	\end{frame}
	
	% =============================================================================
	% SLIDE 4: EL CONFLICTO EMPÍRICO (LAPLACIANO NO NORMALIZADO)
	% =============================================================================
	\begin{frame}{El Giro Dramático: El Colapso del Modelo Ingenuo}
		\begin{columns}[c]
			\begin{column}{0.50\textwidth}
				\textbf{El Operador Laplaciano Estándar:}
				
				```
				
				$$L = D - A, \quad D_{ii} = \sum_j A_{ij}$$
				
				\textbf{Forma Cuadrática (Tensión del Grafo):}
				
				$$x^\top L x = \frac{1}{2} \sum_{i,j} A_{ij}(x_i - x_j)^2$$
				
				```
				\vspace{-0.2cm}
				\begin{alertblock}{El Fracaso del Corte Estándar}
					\begin{itemize}
						\item Minimiza \(\text{cut}(A, B) = \sum_{i \in A, j \in B} A_{ij}\).
						\item \textbf{Patología empírica:} Aísla nodos periféricos con grado mínimo (\(d_i \approx 0\)).
						\item \textbf{Resultado:} Un clúster con el $\sim$98\% del mercado y 3 clústeres triviales de 1 nodo.
						\item \textbf{Adjusted Rand Index (ARI):} $0.021$ (Poder nulo).
					\end{itemize}
				\end{alertblock}
			\end{column}
			
			\begin{column}{0.46\textwidth}
				\begin{center}
					\textbf{\small Topología Resultante: Agrupamiento Trivial}\\
					\vspace{0.2cm}
					% Espacio para gráfico monocromático de fallo empírico
					\begin{tikzpicture}
						\draw[fill=SoftGray, draw=AlertRed, thick] (0,0) rectangle (5.5,3.6);
						\node at (2.75,2.8) {\footnotesize \textbf{Super-Clúster Artificial ($\sim$98\%)}};
						\draw[fill=SlateBlue!40] (2.75,1.5) circle (1.2cm);
						\draw[fill=AlertRed] (0.6,0.6) circle (0.15cm);
						\draw[fill=AlertRed] (4.9,0.8) circle (0.15cm);
						\draw[dashed, red] (1.7,1.1) -- (0.6,0.6);
						\node[below, text width=5cm, align=center] at (2.75,-0.2) {\scriptsize Cortar aristas débiles periféricas es ``barato'' para \(L\), arruinando la macro-partición.};
					\end{tikzpicture}
				\end{center}
			\end{column}
		\end{columns}
		\note{
			[3:15 - 4:55] Miguel:
			Construimos el operador básico: el Laplaciano estándar L = D - A y su forma cuadrática.
			Aquí llegó el gran tropiezo experimental: al ejecutar el algoritmo ingenuo, el modelo colapsó.
			Agrupó al 98% de las acciones en un solo clúster y aisló un par de nodos periféricos.
			¿Por qué? Porque minimizar el corte absoluto es miope: resulta más 'barato' cortar un nodo
			casi aislado con grado bajo que atravesar un puente denso entre dos industrias.
			El álgebra lineal funcionó, pero resolvió la pregunta equivocada.
		}
		
		```
		
	\end{frame}
	
	% =============================================================================
	% SLIDE 5: EL RESCATE MATEMÁTICO (NORMALIZED CUT Y NJW)
	% =============================================================================
	\begin{frame}{El Rescate Matemático: Normalized Cut y el Algoritmo NJW}
		\begin{columns}[c]
			\begin{column}{0.50\textwidth}
				\begin{block}{Penalización por Volumen de Aristas}
					
					$$\text{NCut}(A_1, \dots, A_k) = \sum_{p=1}^k \frac{\text{cut}(A_p, \bar{A}_p)}{\text{vol}(A_p)}$$
					
					```
					donde \(\text{vol}(A_p) = \sum_{i \in A_p} d_i\). Impide grupos diminutos.
					
					```
					
				\end{block}
				\vspace{0.2cm}
				\textbf{Relajación Espectral (Rayleigh-Ritz):}
				
				$$L_{\text{sym}} = D^{-1/2} L D^{-1/2} = I - D^{-1/2} A D^{-1/2}$$
				
				```
				\begin{itemize}
					\item \(L_{\text{sym}}\) es simétrica y semidefinida positiva.
					\item Sus autovalores están acotados: \(0 = \lambda_1 \le \dots \le \lambda_n \le 2\).
				\end{itemize}
				
				```
				
			\end{column}
			
			\begin{column}{0.48\textwidth}
				\begin{block}{El Paso Crítico: Normalización de Filas (NJW)}
					Sea $U = [v_1, \dots, v_k] \in \mathbb{R}^{n \times k}$ con los primeros $k$ autovectores. Se construye $T \in \mathbb{R}^{n \times k}$ mediante:
					
					$$T_{ij} = \frac{U_{ij}}{\left(\sum_{l=1}^k U_{il}^2\right)^{1/2}}$$
					
					```
				\end{block}
				\vspace{0.2cm}
				\textbf{Propiedad Geométrica Fundamental:}
				\begin{itemize}
					\item Cada nodo se proyecta a la esfera unitaria: \(\|T_{i,\cdot}\|_2 = 1\).
					\item Elimina el sesgo de escala por el grado \(d_i\) del nodo en \(K\)-Means.
				\end{itemize}
			\end{column}
		\end{columns}
		\note{
			[4:55 - 7:15] Miguel:
			Para rescatar el agrupamiento recurrimos al Normalized Cut de Shi y Malik.
			En vez de minimizar el corte absoluto, dividimos por el volumen de conexiones.
			Al relajar este problema discreto NP-hard mediante Rayleigh-Ritz, el problema se transforma
			en los autovectores del Laplaciano Normalizado Simétrico L_sym.
			Además, implementamos el paso de Ng, Jordan y Weiss: normalizamos euclidianamente cada fila
			de la matriz U a norma 1, proyectando cada activo sobre una hiperesfera unitaria.
			Esto impide que empresas de alta capitalización o grado extremo distorsionen K-Means.
			Paso la palabra a Jorge para ver qué reveló el espectro.
		}
		
		```
		
	\end{frame}
	
	% =============================================================================
	% SLIDE 6: EL EIGENGAP Y NÚMERO DE COMUNIDADES
	% =============================================================================
	\begin{frame}{La Señal en el Espectro: La Heurística del Eigengap}
		\begin{columns}[c]
			\begin{column}{0.48\textwidth}
				\textbf{Propiedades del Espectro de $L_{\text{sym}}$:}
				\begin{itemize}
					\item $\lambda_1 = 0$ con multiplicidad algebraica $c=1$ (red conexa).
					\item $\lambda_2 > 0$: Conectividad algebraica y cota de Fiedler.
				\end{itemize}
				\vspace{0.2cm}
				\begin{block}{Criterio Objetivo de Elección ($k$)}
					Se busca maximizar el salto espectral relativo:
					
					$$\Delta_k = \lambda_{k+1} - \lambda_k$$
					
					```
					\begin{itemize}
						\item \(\lambda_1, \lambda_2, \lambda_3, \lambda_4\) crecen suavemente.
						\item Quiebre abrupto en \(\lambda_5 - \lambda_4\) (\textbf{Eigengap}).
						\item \textbf{Veredicto matemático:} Fijar \textbf{\(k = 4\)}.
					\end{itemize}
				\end{block}
			\end{column}
			
			\begin{column}{0.50\textwidth}
				\begin{center}
					% Placeholder para la curva del Eigengap obtenida en el Bloque 3
					\begin{tikzpicture}
						\draw[->] (0,0) -- (6,0) node[right] {\scriptsize \(i\)};
						\draw[->] (0,0) -- (0,3.5) node[above] {\scriptsize Magnitud \(\lambda_i\)};
						\draw[blue!70!black, thick, mark=*] plot coordinates {
							(0.5,0.0) (1.0,0.15) (1.5,0.32) (2.0,0.50) (2.5,1.40) (3.0,1.65) (3.5,1.85) (4.0,2.05) (4.5,2.20) (5.0,2.35)
						};
						\draw[red, thick, dashed] (2.0,0.5) -- (2.5,1.4);
						\node[red, right] at (2.2,0.95) {\textbf{\scriptsize Salto Eigengap (\(k=4\))}};
						\node[below] at (2.0,0) {\scriptsize \(\lambda_4\)};
						\node[below] at (2.5,0) {\scriptsize \(\lambda_5\)};
					\end{tikzpicture}\\
					\vspace{0.1cm}
					\small\textit{Espectro del Laplaciano: Los datos piden 4 grupos sin forzar la hipótesis.}
				\end{center}
			\end{column}
		\end{columns}
		\note{
			[7:15 - 8:30] Jorge:
			Gracias, Miguel. Al descomponer espectralmente L_sym, ordenamos sus autovalores.
			$\Lambda_1$ es 0 y tiene multiplicidad 1 porque nuestra red es completamente conexa.
			Para fijar k de manera rigurosa sin sesgos usamos la heurística del Eigengap,
			buscando el salto máximo entre autovalores consecutivos.
			Como observan en la curva, los primeros 4 autovalores son pequeños y crecen suavemente,
			pero tras $\lambda_4$ se produce un quiebre pronunciado hacia $\lambda_5$.
			El espectro reveló de forma intrínseca que el grafo se desacopla naturalmente en 4 comunidades.
		}
		
		```
		
	\end{frame}
	
	% =============================================================================
	% SLIDE 7: EL EMBEDDING ESPECTRAL
	% =============================================================================
	\begin{frame}{El Espacio Latente: Mapeo Previo al Clustering}
		\begin{columns}[c]
			\begin{column}{0.48\textwidth}
				\textbf{De Nodos a Coordenadas Geométricas:}
				\begin{itemize}
					\item Cada fila de $T$ ubica al activo en $\mathbb{R}^4$.
					\item Eje X: \textbf{Vector de Fiedler ($v_2$)} $\to$ Bipartición primaria de la red.
					\item Eje Y: \textbf{Autovector 3 ($v_3$)} $\to$ Segunda dirección de máxima varianza espectral.
				\end{itemize}
				\vspace{0.2cm}
				\begin{block}{Desacoplamiento sin Supervisión}
					Sin ejecutar $K$-Means aún, los sectores se auto-organizan en cuadrantes separados:
					\begin{itemize}
						\item \textbf{Energía:} Rayo radial aislado.
						\item \textbf{Finanzas:} Nube compacta superior.
						\item \textbf{Tecnología y Salud:} Separadas ortogonalmente.
					\end{itemize}
				\end{block}
			\end{column}
			
			```
			\begin{column}{0.50\textwidth}
				\begin{center}
					% Espacio para la gráfica de dispersión 2D del Embedding
					\begin{tikzpicture}[scale=0.9]
						\draw[->, gray] (-2.5,0) -- (2.5,0) node[right, black] {\scriptsize \(v_2\) (Fiedler)};
						\draw[->, gray] (0,-2.2) -- (0,2.2) node[above, black] {\scriptsize \(v_3\)};
						% Nubes de puntos representativas
						\draw[fill=orange!80, opacity=0.8] plot coordinates {(1.5,0.2) (1.8,0.4) (1.6,-0.1) (2.0,0.1) (1.7,0.3)};
						\node[right, orange!90!black] at (1.8,-0.3) {\scriptsize Energía};
						\draw[fill=blue!70, opacity=0.8] plot coordinates {(0.1,1.4) (-0.2,1.6) (0.3,1.7) (-0.1,1.9) (0.2,1.3)};
						\node[above, blue!80!black] at (0.2,1.9) {\scriptsize Finanzas};
						\draw[fill=red!70, opacity=0.8] plot coordinates {(-1.4,0.3) (-1.7,0.5) (-1.3,0.1) (-1.8,0.2)};
						\node[left, red!80!black] at (-1.8,0.3) {\scriptsize Tecnología};
						\draw[fill=green!70!black, opacity=0.8] plot coordinates {(-0.5,-1.2) (-0.8,-1.5) (-0.3,-1.4) (-0.6,-1.7)};
						\node[below, green!70!black] at (-0.5,-1.7) {\scriptsize Salud};
					\end{tikzpicture}\\
					\vspace{0.1cm}
					\small\textit{Proyección Espectral: La distancia euclidiana refleja disimilitud topológica.}
				\end{center}
			\end{column}
		\end{columns}
		\note{
			[8:30 - 9:55] Jorge:
			Antes de ejecutar K-Means, visualizamos el embedding espectral proyectando las acciones
			en el plano del vector de Fiedler v_2 y el autovector v_3.
			Aquí los nodos pasan de ser vértices a coordenadas euclidianas en R^4.
			Observen el fenómeno: sin haber corrido clustering aún, el Laplaciano ya separó los sectores.
			Las energéticas forman un rayo hacia la derecha; las finanzas están compactas arriba;
			y salud con tecnología ocupan cuadrantes opuestos.
			K-Means simplemente formaliza las fronteras que el espectro ya resolvió algebraicamente.
		}
		
		```
		
	\end{frame}
	
	% =============================================================================
	% SLIDE 8: LA RED RECONSTRUIDA
	% =============================================================================
	\begin{frame}{Topología de la Red: Comunidades Bursátiles Reconstruidas}
		\begin{columns}[c]
			\begin{column}{0.48\textwidth}
				\textbf{Reconstrucción de la Red Compleja:}
				\begin{itemize}
					\item Visualización con algoritmo de fuerzas \textit{Spring Layout} de NetworkX.
					\item Cuatro colores asignados por $K$-Means sobre la representación espectral $T$.
				\end{itemize}
				\vspace{0.2cm}
				\textbf{Validación de Macro-Sectores:}
				\begin{itemize}
					\item \textbf{Energía:} Núcleo ultra-cohesivo (ej. \texttt{XOM}, \texttt{CVX}, \texttt{SLB}).
					\item \textbf{Banca e Inversión:} Bloque compacto (\texttt{JPM}, \texttt{GS}, \texttt{BAC}).
					\item \textbf{Farmacéutica y Salud:} Sin mezcla con finanzas (\texttt{PFE}, \texttt{JNJ}, \texttt{UNH}).
					\item \textbf{Tecnología y Software:} Alta conectividad intra-cluster (\texttt{AAPL}, \texttt{MSFT}, \texttt{NVDA}).
				\end{itemize}
			\end{column}
			
			```
			\begin{column}{0.50\textwidth}
				\begin{center}
					% Espacio para la figura exportada por NetworkX
					\begin{tikzpicture}
						\draw[fill=SoftGray, draw=NavyDark, thick] (0,0) rectangle (6,3.8);
						\node at (3,2) {\textit{\small [Gráfico NetworkX: 141 Nodos / 4 Comunidades]}};
						\node[below] at (3,1.4) {\scriptsize \textbf{Rojo}: Tech | \textbf{Azul}: Finanzas | \textbf{Verde}: Salud | \textbf{Naranja}: Energía};
					\end{tikzpicture}\\
					\vspace{0.1cm}
					\small\textit{Grafo S\&P 500: Separación limpia de la densidad estructural de covarianza.}
				\end{center}
			\end{column}
		\end{columns}
		\note{
			[9:55 - 11:15] Jorge:
			Al proyectar las etiquetas sobre el grafo topológico, el resultado fue contundente:
			el algoritmo reconstruyó a ciegas las comunidades sectoriales del mercado.
			El bloque de Energía quedó encapsulado con aristas densas internas;
			el clúster Financiero agrupó a la banca y aseguradoras aislado del resto;
			y Tecnología con Salud se consolidaron como macro-comunidades estables.
			Aquel colapso del clúster único que vimos con el Laplaciano estándar quedó superado.
		}
		
		```
		
	\end{frame}
	
% =============================================================================
% SLIDE 9: VALIDACIÓN CUANTITATIVA Y DISEÑO EXPERIMENTAL
% =============================================================================
\begin{frame}{Diseño Experimental y Métricas de Partición}
	\begin{columns}[T]
		\begin{column}{0.55\textwidth}
			\textbf{Comparación Sistemática de Modelos:}
			\vspace{0.2cm}
			
			% TABLA 1 CORREGIDA
			\begin{table}[h]
				\centering
				\scriptsize
				\begin{tabular}{l c c c c}
					\toprule
					\textbf{Configuración} & \textbf{ARI} \(\uparrow\) & \textbf{NMI} \(\uparrow\) & \textbf{Silueta} & \textbf{Q} \\
					\midrule
					\(L\) No Normal. & 0,021 & 0,045 & -0,12 & 0,031 \\
					Shi-Malik (\(L_{\text{rw}}\)) & 0,812 & 0,774 & 0,46 & 0,442 \\
					\textbf{NJW (\(L_{\text{sym}}\))} & \textbf{0,835} & \textbf{0,791} & \textbf{0,49} & \textbf{0,468} \\
					Scikit-Learn  & 0,835 & 0,791 & 0,49 & 0,468 \\
					\bottomrule
				\end{tabular}
			\end{table}
			
			\vspace{0.1cm}
			\begin{itemize}
				\item \textbf{Decisión 1 (\(L\) vs \(L_{\text{sym}}\)):} El salto en ARI a $0,83$ demuestra la necesidad de penalizar el volumen.
				\item \textbf{Decisión 2 (Benchmark):} Coincidencia exacta con la librería Scikit-Learn.
			\end{itemize}
		\end{column}
		
		\begin{column}{0.43\textwidth}
			\textbf{Matriz de Confusión:}
			\vspace{0.1cm}
			
			% TABLA 2 CORREGIDA
			\begin{table}[h]
				\centering
				\tiny
				\begin{tabular}{l c c c c}
					\toprule
					\textbf{Sector} & \textbf{C1} & \textbf{C2} & \textbf{C3} & \textbf{C4} \\
					\midrule
					Finanzas   & \textbf{33} & 2  & 0  & 0  \\
					Tecnología & 1  & \textbf{34} & 1  & 0  \\
					Salud      & 0  & 2  & \textbf{34} & 0  \\
					Energía    & 0  & 2  & 0  & \textbf{32} \\
					\bottomrule
				\end{tabular}
			\end{table}
			
			\vspace{0.1cm}
			\small\textit{Altísima pureza diagonal; los cruces corresponden a activos con dinámicas híbridas.}
		\end{column}
	\end{columns}
	\note{
		[11:15 - 12:35] Jorge:
		Medimos el desempeño cuantitativo del diseño experimental.
		Con el Laplaciano no normalizado, el Adjusted Rand Index fue nulo.
		Al migrar a L\_sym, el ARI saltó a 0.835 y el NMI a casi 0.80, respaldando la coincidencia
		con los sectores reales.
		Además, contrastamos contra Scikit-Learn obteniendo resultados idénticos,
		lo que garantiza la reproducibilidad de nuestro código.
	}
\end{frame}
	
	% =============================================================================
	% SLIDE 10: NODOS FRONTERIZOS E INSIGHTS FINANCIEROS
	% =============================================================================
	\begin{frame}{Nodos Fronterizos: Cuando el Riesgo Desafía la Etiqueta}
		\begin{columns}[c]
			\begin{column}{0.50\textwidth}
				\textbf{El Hallazgo en las Fronteras Espectrales:}
				\begin{itemize}
					\item Las clasificaciones GICS ubican a \texttt{FSLR} (First Solar) y \texttt{ENPH} (Enphase Energy) en Energía o Tecnología.
					\item El clustering espectral las situó en la \textbf{frontera euclidiana} entre ambos mundos.
				\end{itemize}
				\vspace{0.2cm}
				\begin{block}{Explicación de Riesgo Real}
					\begin{itemize}
						\item \textbf{Desacople del Petróleo:} Su covarianza no depende del barril de crudo tradicional (WTI/Brent).
						\item \textbf{Sensibilidad Tecnológica:} Dependen de tasas de descuento, semiconductores y subsidios de innovación limpia.
					\end{itemize}
				\end{block}
			\end{column}
			
			```
			\begin{column}{0.48\textwidth}
				\begin{alertblock}{Conclusión del Analista Cuantitativo}
					El clustering espectral no cometió un error de clasificación: \\
					\vspace{0.1cm}
					\textbf{Descubrió que el riesgo sistémico de la transición verde desacopla a estas empresas de su etiqueta legal histórica.}
				\end{alertblock}
				\vspace{0.3cm}
				\centering
				\begin{tikzpicture}
					\node[draw, circle, fill=orange!40, minimum size=1cm] (E) at (0,0) {\scriptsize \textbf{Energía}};
					\node[draw, circle, fill=red!30, minimum size=1cm] (T) at (3.5,0) {\scriptsize \textbf{Tech}};
					\node[draw, rectangle, rounded corners, fill=yellow!60, draw=red, thick] (S) at (1.75,0) {\tiny \texttt{FSLR / ENPH}};
					\draw[<->, thick] (E) -- (S);
					\draw[<->, thick] (S) -- (T);
				\end{tikzpicture}
			\end{column}
		\end{columns}
		\note{
			[12:35 - 13:45] Jorge:
			El hallazgo más valioso provino de los 'errores' de clasificación: los nodos fronterizos.
			Empresas como First Solar y Enphase están catalogadas en Energía o Tecnología tradicional,
			pero el algoritmo las situó en la frontera periférica.
			¿Por qué? Porque son empresas de energía solar: sus retornos no covarían con el barril
			de crudo de Exxon o Chevron, sino con el ciclo de tasas de interés y componentes tecnológicos.
			El modelo no falló: detectó que el riesgo financiero real desafía la etiqueta burocrática.
		}
		
		```
		
	\end{frame}
	
	% =============================================================================
	% SLIDE 11: LIMITACIONES Y EXTENSIONES
	% =============================================================================
	\begin{frame}{Limitaciones y Trabajo Futuro}
		\begin{columns}[t]
			\begin{column}{0.48\textwidth}
				\begin{block}{Limitaciones del Modelo Actual}
					\begin{enumerate}
						\item \textbf{Correlación Estática:} Asume relaciones lineales invariantes en 10 años, ignorando quiebres en crisis.
						\item \textbf{Sensibilidad al Umbral:} El filtro de corte ($\rho \ge 0.40$) requiere calibración empírica.
						\item \textbf{Convexidad en $K$-Means:} Asume nubes hiper-esféricas tras la normalización.
					\end{enumerate}
				\end{block}
			\end{column}
			
			```
			\begin{column}{0.48\textwidth}
				\begin{block}{Extensiones Propuestas}
					\begin{enumerate}
						\item \textbf{Redes Temporales:} Ventanas móviles de correlación para estudiar la evolución del espectro durante pánicos de liquidez.
						\item \textbf{Comunidades Superpuestas:} Algoritmos de factorización no negativa (NMF) sobre el Laplaciano.
						\item \textbf{Graph Neural Networks (GNNs):} Integrar la topología del grafo con balances fundamentales contables.
					\end{enumerate}
				\end{block}
			\end{column}
		\end{columns}
		\note{
			[13:45 - 14:20] Jorge:
			Como limitaciones principales, la correlación de Pearson asume relaciones lineales estáticas
			y el umbral de 0.40 es sensible.
			Como extensiones, proponemos construir redes dinámicas con ventanas móviles para observar
			cómo se reconfigura el Laplaciano en crisis; comparar con optimización de modularidad vía Louvain;
			o implementar Graph Neural Networks para combinar topología con métricas fundamentales.
		}
		
		```
		
	\end{frame}
	
% =============================================================================
% SLIDE 12: CONCLUSIONES Y DEFENSA
% =============================================================================
\begin{frame}{Conclusiones y Cierre}
	\begin{block}{Síntesis del Trabajo}
		\begin{itemize}
			\item \textbf{Conexión Matemática:} El clustering espectral transforma la partición combinatoria NP-dura en una relajación continua elegante resuelta por autovectores.
			\item \textbf{Relevancia del Normalized Cut:} El Laplaciano normalizado simétrico (\(L_{\text{sym}}\)) con proyección de filas es indispensable en grafos reales no balanceados.
			\item \textbf{Aporte a Finanzas Cuantitativas:} El espectro extrae la macro-estructura latente del S\&P 500 y señala empresas en transición antes que la taxonomía oficial.
		\end{itemize}
	\end{block}
	\vspace{0.4cm}
	\centering
	\textbf{\Large Muchas gracias}\\
	\vspace{0.2cm}
	\textit{Quedamos a disposición de la comisión para la ronda de preguntas.}\\
	\vspace{0.3cm}
	\footnotesize\textbf{Repositorio y Código Reproducible disponible en la entrega.}
	\note{
		[14:20 - 14:40] Miguel:
		Para concluir: el clustering espectral demostró ser un puente formidable entre el álgebra
		lineal y la ciencia de datos en redes complejas.
		Transformó un problema combinatorio intratable en un cálculo de autovalores estable,
		filtrando el ruido del mercado y revelando la verdadera estructura del riesgo.
		Muchas gracias. Quedamos a su entera disposición para sus preguntas.
	}
\end{frame}
	
\end{document}
